% =========================================================================
\section{La ley de Gauss}\label{sec:gauss}
% =========================================================================
Comenzamos a estudiar las propiedades del campo electrostático. La naturaleza de \emph{cuadrado inverso} de la ley de
Coulomb permite conocer el flujo de $\vE$ a través de cualquier superficie cerrada y, de ahí, su divergencia. La herramienta
geométrica que lo hace posible es el ángulo sólido.

% -------------------------------------------------------------------------
\subsection{El ángulo sólido}\label{sec:angulo-solido}
% -------------------------------------------------------------------------
La razón entre las áreas totales de dos esferas concéntricas y los cuadrados de sus radios es constante:
$4\pi r_1^{2}/r_1^{2}=4\pi r_2^{2}/r_2^{2}=4\pi$. Lo mismo ocurre con áreas elementales vistas desde el centro
común $O$ (\figpanel{fig:angulo}{a}): si dos elementos $\dif S_1$ y $\dif S_2$ de esferas de radios $r_1$ y $r_2$ están
recortados por el mismo cono con vértice en $O$, entonces $\dif S_1/r_1^{2}=\dif S_2/r_2^{2}$. Esa constante es el
ángulo sólido, la versión tridimensional del ángulo en radianes (que es la razón entre un arco y su radio).

\begin{definicion}[Ángulo sólido]\label{def:angulo-solido}
El \emph{ángulo sólido elemental} bajo el que se ve desde un punto $O$ un elemento de superficie esférica
$\dif S$ de radio $r$ centrado en $O$ es
\begin{equation}\label{eq:angulo-esfera}
  \dOm=\frac{\dif S_1}{r_1^{2}}=\frac{\dif S_2}{r_2^{2}}
\end{equation}
(\emph{estereorradianes}, sr). Con $r=1$, es el área de la porción de esfera unidad bajo la cual se ve el elemento.
Para un elemento $\dS=\un\,\dsc$ \emph{cualquiera}, situado en un punto cuyo vector de posición relativa a $O$ es
$\vc{R}$, se proyecta sobre la esfera de radio $R$ centrada en $O$ (\figpanel{fig:angulo}{b}): el área proyectada
es $\dif S'=\dsc\cos\theta$, con $\theta$ el ángulo entre $\dS$ y $\vc{R}$, y
\begin{equation}\label{eq:angulo-general}
  \dOm=\frac{\dif S'}{R^{2}}=\frac{\dsc\cos\theta}{R^{2}}=\frac{\dS\cdot\vc{R}}{R^{3}}=\frac{\dS\cdot\uR}{R^{2}}.
\end{equation}
El signo de $\dOm$ es el de $\cos\theta$: positivo si $\vc{R}$ atraviesa el elemento «de dentro hacia fuera»
(según $\un$) y negativo en caso contrario.
\end{definicion}

% Figura: ángulo sólido. (a) esferas concéntricas, (b) elemento oblicuo y su proyección,
% (c) punto exterior: dos elementos opuestos recortados por el mismo cono.
\begin{figure}[ht]
\centering
\tikzset{
  parche/.style={draw=azulprof, line width=0.6pt, fill=azulclaro!70, fill opacity=0.9},
  parche2/.style={draw=acento, line width=0.6pt, fill=acento!25, fill opacity=0.9},
  cono/.style={azulacero, line width=0.45pt},
  vn/.style={-{Stealth[length=5pt,width=4.5pt]}, line width=1.1pt}}
\begin{minipage}[t]{0.50\textwidth}\centering
\begin{tikzpicture}[baseline=(current bounding box.north), x=1cm, y=1cm]
  % (a) dos esferas concéntricas
  \shade[ball color=azulclaro!55, opacity=0.30] (0,0) circle (2.3);
  \draw[azulprof, line width=0.8pt] (0,0) circle (2.3);
  \shade[ball color=azulclaro!70, opacity=0.35] (0,0) circle (1.15);
  \draw[azulacero, line width=0.8pt] (0,0) circle (1.15);
  % ecuadores (círculos máximos en perspectiva)
  \draw[acento, line width=0.8pt] (-2.3,0) arc (180:360:{2.3} and 0.5);
  \draw[acento, line width=0.5pt, densely dashed] (-2.3,0) arc (180:0:{2.3} and 0.5);
  \draw[acento, line width=0.8pt] (-1.15,0) arc (180:360:{1.15} and 0.25);
  \draw[acento, line width=0.5pt, densely dashed] (-1.15,0) arc (180:0:{1.15} and 0.25);
  % cono hacia los dos parches (dirección 52 grados)
  \def\dirr{48}
  \def\ra{0.86}\def\rb{1.72}
  \draw[cono] (0,0) -- ($(\dirr+12:\rb)$) (0,0) -- ($(\dirr-12:\rb)$);
  \begin{scope}[shift={(\dirr:\ra)}, rotate=\dirr+90]
    \path[parche] (0,0) ellipse (0.19 and 0.075);
  \end{scope}
  \begin{scope}[shift={(\dirr:\rb)}, rotate=\dirr+90]
    \path[parche] (0,0) ellipse (0.38 and 0.15);
  \end{scope}
  \fill[azulprof] (0,0) circle (1.6pt);
  \node[font=\small, anchor=north east] at (-0.05,-0.12) {$O$};
  \node[font=\small, text=azulprof, anchor=east] at ($(\dirr:\ra)+(-0.22,0.1)$) {$\dif S_1$};
  \node[font=\small, text=azulprof, anchor=south west] at ($(\dirr:\rb)+(0.2,0.12)$) {$\dif S_2$};
  % radios sobre el plano ecuatorial
  \draw[vn, verdeteo, line width=0.9pt] (0,0) -- (194:1.08) node[pos=0.55, above, font=\scriptsize, text=verdeteo] {$r_1$};
  \draw[vn, verdeteo, line width=0.9pt] (0,0) -- (-12:2.2) node[pos=0.72, above, font=\scriptsize, text=verdeteo] {$r_2$};
\end{tikzpicture}\\[2pt]{\small\scshape (a)}
\end{minipage}\hfill
\begin{minipage}[t]{0.49\textwidth}\centering
\begin{tikzpicture}[baseline=(current bounding box.north), x=1cm, y=1cm]
  % (b) superficie cerrada S, esfera unidad y esfera S' de radio R
  \draw[azulprof, line width=0.9pt, fill=azulfondo, fill opacity=0.6]
    plot[smooth cycle, tension=0.75] coordinates {(-0.4:2.7) (30:3.0) (68:3.4) (120:3.0) (172:2.6) (220:3.0) (268:2.6) (318:3.1)};
  \draw[azulacero, line width=0.5pt] (0,0) circle (0.9);
  \draw[azulacero, line width=0.9pt] (6:2.9) arc (6:70:2.9);
  \draw[azulacero, line width=0.9pt, densely dashed] (70:2.9) arc (70:92:2.9);
  \draw[azulacero, line width=0.9pt, densely dashed] (6:2.9) arc (6:-26:2.9);
  \fill[azulprof] (0,0) circle (1.6pt); \node[font=\small, anchor=north east] at (-0.05,-0.1) {$O$};
  \node[font=\scriptsize, text=azulprof] at (0.0,0.42) {$1$};
  \node[font=\small, text=azulprof] at (150:2.2) {$S$};
  \node[font=\small, text=azulacero] at (-20:2.45) {$S'$};
  % cono
  \def\dirr{30}
  \draw[cono] (0,0) -- ($(\dirr+9:2.9)$) (0,0) -- ($(\dirr-9:2.9)$);
  \begin{scope}[shift={(\dirr:0.9)}, rotate=\dirr+90] \path[parche] (0,0) ellipse (0.14 and 0.05); \end{scope}
  % dS' (proyección sobre S') y dS (inclinado un ángulo theta)
  \begin{scope}[shift={(\dirr:2.9)}, rotate=\dirr+90] \path[parche] (0,0) ellipse (0.46 and 0.15); \end{scope}
  \begin{scope}[shift={(\dirr:2.9)}, rotate=\dirr+90+30] \path[parche2] (0,0) ellipse (0.52 and 0.16); \end{scope}
  \draw[vn, azulacero] (\dirr:2.9) -- ($(\dirr:2.9)+(\dirr:0.95)$) node[right, font=\small, text=azulacero] {$\dif\vc{S}'$};
  \draw[vn, acento] (\dirr:2.9) -- ($(\dirr:2.9)+(\dirr+30:0.95)$) node[left=1pt, font=\small, text=acento] {$\dif\vc{S}$};
  \draw[verdeteo, line width=0.5pt] ($(\dirr:2.9)+(\dirr:0.55)$) arc (\dirr:\dirr+30:0.55);
  \node[font=\scriptsize, text=verdeteo] at ($(\dirr+15:2.9)+(\dirr+15:0.85)$) {$\theta$};
  \draw[vn, azulprof, line width=0.9pt] (0,0) -- ($(\dirr:2.9)-(\dirr:0.2)$) node[pos=0.55, below right=-1pt, font=\small, text=azulprof] {$\vc{R}$};
\end{tikzpicture}\\[2pt]{\small\scshape (b)}
\end{minipage}\\[10pt]
\begin{minipage}[t]{0.8\textwidth}\centering
\begin{tikzpicture}[baseline=(current bounding box.north), x=1cm, y=1cm]
  % (c) O fuera de S: el cono recorta dos elementos en lados opuestos
  \def\ang{7.5}
  \fill[azulprof] (0,0) circle (1.6pt); \node[font=\small, anchor=east] at (-0.1,0) {$O$};
  \draw[azulprof, line width=0.9pt, fill=azulfondo, fill opacity=0.6]
    plot[smooth cycle, tension=0.8] coordinates {(4.2,0.7) (5.0,1.25) (6.2,1.3) (7.3,0.95) (7.9,-0.1) (7.2,-1.2) (5.9,-1.3) (4.7,-0.9)};
  \draw[cono] (0,0) -- (7.55,{7.55*tan(\ang)}) (0,0) -- (7.55,{-7.55*tan(\ang)});
  \draw[azulprof, densely dashed, line width=0.6pt] (0,0) -- (6.9,0);
  % parche de la esfera unidad
  \draw[azulacero, line width=0.5pt] (0.9,0.6) arc (36:-36:1.0);
  \begin{scope}[shift={(1.0,0)}] \path[parche] (0,0) ellipse (0.045 and 0.135); \end{scope}
  \node[font=\scriptsize, text=azulprof, anchor=south] at (0.55,0.07) {$1$};
  % elemento cercano dS2 (entrada) y lejano dS1 (salida)
  \begin{scope}[shift={(4.85,0)}, rotate=20] \path[parche2] (0,0) ellipse (0.14 and 0.46); \end{scope}
  \begin{scope}[shift={(7.35,0)}, rotate=-10] \path[parche] (0,0) ellipse (0.18 and 0.62); \end{scope}
  \draw[vn, acento] (4.85,0) -- ++(150:0.95) node[above left=-1pt, font=\small, text=acento] {$\dif\vc{S}_2$};
  \draw[vn, azulacero] (7.35,0) -- ++(25:0.95) node[above right, font=\small, text=azulacero] {$\dif\vc{S}_1$};
  \draw[verdeteo, line width=0.5pt] (5.3,0) arc (0:150:0.45);
  \node[font=\scriptsize, text=verdeteo] at (4.55,0.55) {$\theta_2$};
  \draw[verdeteo, line width=0.5pt] (7.75,0) arc (0:25:0.4);
  \node[font=\scriptsize, text=verdeteo, anchor=west] at (7.78,0.12) {$\theta_1$};
  \node[font=\scriptsize, text=azulprof, anchor=north] at (2.6,-0.05) {$\vc{R}_2$};
  \node[font=\scriptsize, text=azulprof, anchor=north] at (6.15,-0.05) {$\vc{R}_1$};
  \node[font=\small, text=azulprof] at (6.3,0.85) {$S$};
\end{tikzpicture}\\[2pt]{\small\scshape (c)}
\end{minipage}
\caption{Ángulo sólido. (a)~Dos elementos de esferas concéntricas, $\dif S_1$ y $\dif S_2$, recortados por el mismo cono de vértice $O$: $\dif S_1/r_1^{2}=\dif S_2/r_2^{2}=\dOm$.
(b)~Para un elemento $\dS$ cualquiera de una superficie cerrada $S$ que rodea a $O$ basta proyectarlo sobre la esfera $S'$ de radio $R$: $\dif S'=\dsc\cos\theta$ y $\dOm=\dif S'/R^{2}$.
(c)~Si $O$ queda fuera de $S$, el mismo cono recorta dos elementos en lados opuestos de $S$; el de entrada ($\theta_2>\pi/2$) aporta $-\dOm$ y el de salida ($\theta_1<\pi/2$) aporta $+\dOm$, y se anulan.}
\label{fig:angulo}
\end{figure}


\begin{lema}[Ángulo sólido total de una superficie cerrada]\label{lema:omega-4pi}
Sea $S$ una superficie cerrada, orientada con la normal exterior. Entonces
\begin{equation}\label{eq:omega}
  \oint_{S}\dOm=\oint_{S}\frac{\dS\cdot\uR}{R^{2}}=\begin{cases}4\pi,&\text{si $O$ está en el interior de $S$,}\\[2pt]0,&\text{si $O$ está en el exterior de $S$.}\end{cases}
\end{equation}
\end{lema}
\begin{proof}
\emph{(a) $O$ en el interior.} Cada semirrecta que sale de $O$ corta a $S$ un número impar de veces, y la
normal exterior forma ángulos alternadamente agudos (salida) y obtusos (entrada) con la dirección de la semirrecta:
$+\dOm_0,-\dOm_0,+\dOm_0,\dots$, donde $\dOm_0$ es el ángulo sólido del cono. Como el último cruce es una
salida, la contribución neta de la dirección es $+\dOm_0$, y sumando sobre todas las direcciones se obtiene el ángulo sólido
total del espacio. Para una superficie estrellada respecto de $O$ (un único cruce por semirrecta) basta proyectar sobre una esfera
$S'$ de radio $R$ centrada en $O$, en la que $R$ es constante:
\[
  \Omega=\oint_{S}\dOm=\oint_{S'}\frac{\dif S'}{R^{2}}=\frac{1}{R^{2}}\oint_{S'}\dif S'=\frac{4\pi R^{2}}{R^{2}}=4\pi .
\]
\emph{(b) $O$ en el exterior.} Ahora cada semirrecta que corta a $S$ lo hace un número \emph{par} de veces, con entradas y
salidas alternadas (\figpanel{fig:angulo}{c}). Para un elemento de entrada, $\theta_2>\pi/2$ y por tanto
$\dS_2\cdot\vc{R}_2/R_2^{3}=\dsc_2\cos\theta_2/R_2^{2}=-\dOm_0$, mientras que para el de salida
$\dS_1\cdot\vc{R}_1/R_1^{3}=\dsc_1\cos\theta_1/R_1^{2}=+\dOm_0$. Se anulan por parejas y la integral es nula.
\end{proof}

% -------------------------------------------------------------------------
\subsection{Ley de Gauss en forma integral}\label{sec:gauss-integral}
% -------------------------------------------------------------------------
Calculemos ahora el flujo del campo eléctrico $\flujo=\oint_{S}\vE\cdot\dS$ en distintos supuestos. Observemos de entrada que
para una carga puntual en $\vc{r}\,'$, con $\vc{R}=\vc{r}-\vc{r}\,'$, el integrando contiene exactamente el ángulo sólido:
\[
  \vE\cdot\dS=\frac{\qp}{4\pi\eo}\,\frac{\uR\cdot\dS}{R^{2}}=\frac{\qp}{4\pi\eo}\,\dOm .
\]

\begin{teorema}[Ley de Gauss, forma integral]\label{teo:gauss-integral}
El flujo del campo eléctrico a través de una superficie cerrada cualquiera $S$ es igual a la carga
encerrada por ella dividida entre $\eo$:
\begin{equation}\label{eq:gauss-int}
  \boxed{\;\flujo=\oiint_{S}\vE\cdot\dS=\frac{\Qenc}{\eo}\;}
\end{equation}
con independencia de que la carga sea una distribución discreta o continua. Si hay una distribución volúmica,
$\Qenc=\int_{V}\rho\,\dV$, con $V$ el volumen contenido en $S$.
\end{teorema}
\begin{proof}
\emph{Carga puntual fuera de $S$.} Usando el \cref{lema:omega-4pi}(b),
\[
  \flujo=\oiint_{S}\frac{\qp}{4\pi\eo}\frac{\uR\cdot\dS}{R^{2}}=\frac{\qp}{4\pi\eo}\oiint_{S}\dOm=\frac{\qp}{4\pi\eo}\,(\Omega-\Omega)=0 .
\]
\emph{Carga puntual dentro de $S$.} Por el \cref{lema:omega-4pi}(a),
\[
  \flujo=\frac{\qp}{4\pi\eo}\oiint_{S}\dOm=\frac{\qp}{4\pi\eo}\,4\pi=\frac{\qp}{\eo}.
\]
\emph{Sistema de cargas puntuales.} Por superposición, el flujo es la suma de los flujos de cada carga; solo las
interiores aportan:
\[
  \flujo=\sum_{\text{dentro}}\frac{\qp_i}{\eo}+\sum_{\text{fuera}}0=\frac{\Qenc}{\eo}.
\]
\emph{Distribución continua.} Se aplica el resultado anterior a cada elemento $\dif q'=\rho\,\dVp$ y se integra:
$\flujo=\frac{1}{\eo}\int_{V}\rho\,\dV$.
\end{proof}

\begin{observacion}[Alcance de la ley de Gauss]\label{obs:alcance-gauss}
\begin{enumerate}[label=(\alph*)]
  \item El flujo de $\vE$ a través de una superficie cerrada depende \emph{solo} de la carga encerrada.
  \item El valor del campo puede variar de un punto a otro de la superficie, ser nulo o no, aunque el flujo sea nulo: en un
        dipolo, el flujo a través de una superficie que lo rodee es nulo, pero el campo no se anula en ningún punto de ella.
  \item Las cargas exteriores no afectan al flujo, pero sí al valor del campo en la superficie.
  \item La posición de las cargas interiores no afecta al flujo, pero sí al valor del campo.
\end{enumerate}
\end{observacion}

% Figura: flujo de una carga puntual dentro y fuera de una superficie cerrada.
\begin{figure}[ht]
\centering
\begin{minipage}[t]{0.48\textwidth}\centering
\begin{tikzpicture}[baseline=(current bounding box.north), scale=0.95]
  \draw[azulprof, line width=0.8pt, fill=azulfondo, fill opacity=0.55] plot[smooth cycle, tension=0.8] coordinates {(1.9,0.2) (1.3,1.4) (-0.2,1.7) (-1.6,1.1) (-1.9,-0.3) (-0.9,-1.5) (0.6,-1.3)};
  \node[font=\small, text=azulprof] at (1.35,-1.05) {$S$};
  \foreach \a in {0,30,...,330}{\draw[-{Stealth[length=4pt]}, cm1, line width=0.6pt] (0.2,0.0) ++(\a:0.15) -- ++(\a:2.5);}
  \node[cargapos] at (0.2,0.0) {$+$};
  \node[font=\small, anchor=south west] at (0.3,0.1) {$\qp$};
\end{tikzpicture}\\[2pt]{\small\scshape (a)}
\end{minipage}\hfill
\begin{minipage}[t]{0.48\textwidth}\centering
\begin{tikzpicture}[baseline=(current bounding box.north), scale=0.95]
  \draw[azulprof, line width=0.8pt, fill=azulfondo, fill opacity=0.55, name path=blob] plot[smooth cycle, tension=0.8] coordinates {(1.9,0.2) (1.3,1.4) (-0.2,1.7) (-1.6,1.1) (-1.9,-0.3) (-0.9,-1.5) (0.6,-1.3)};
  \node[font=\small, text=azulprof] at (1.35,-1.05) {$S$};
  \coordinate (Q) at (-3.4,0.1);
  \foreach \a in {-34,-22,-10,2,14,26}{
    \path[name path=r\a] (Q) -- ++(\a:7.0);
    \path[name intersections={of=blob and r\a, name=I\a, total=\t}];
    \draw[-{Stealth[length=4pt]}, cm1, line width=0.6pt] (Q) ++(\a:0.15) -- ++(\a:6.0);}
  \foreach \a in {-10,2,14}{
    \path[name path=r\a] (Q) -- ++(\a:7.0);
    \path[name intersections={of=blob and r\a, sort by=r\a, name=J\a}];
    \draw[azulprof, fill=white, line width=0.5pt] (J\a-1) circle (1.8pt);
    \fill[azulprof] (J\a-2) circle (1.8pt);}
  \node[cargapos] at (Q) {$+$};
  \node[font=\small, anchor=south east] at ($(Q)+(-0.1,0.1)$) {$\qp$};
\end{tikzpicture}\\[2pt]{\small\scshape (b)}
\end{minipage}
\caption{Flujo de una carga puntual. (a)~Carga interior: todas las líneas de campo que parten de $\qp$
atraviesan $S$ hacia fuera y el flujo vale $\qp/\eo$ (equivale a $4\pi$ de ángulo sólido). (b)~Carga exterior:
cada línea que entra por $S$ ($\circ$) sale después por otro punto ($\bullet$); el flujo de entrada compensa
el de salida y el flujo neto es nulo, aunque el campo no lo sea en $S$.}
\label{fig:flujo-gauss}
\end{figure}


% -------------------------------------------------------------------------
\subsection{Ley de Gauss en forma diferencial}\label{sec:gauss-diferencial}
% -------------------------------------------------------------------------
La forma integral es una afirmación sobre regiones extensas; la forma diferencial es su traducción local, punto a punto. Para
pasar de una a otra usamos el teorema de la divergencia del Tema~0 (\cref{eq:div-teorema} más abajo).

\begin{teorema}[Teorema de la divergencia o de Gauss~\cite{TEMA0}]\label{teo:divergencia}
Si $\vc{A}$ es un campo vectorial de clase $C^{1}$ en un entorno de un sólido acotado $V$ de borde $S=\partial V$ orientado con
la normal exterior,
\begin{equation}\label{eq:div-teorema}
  \iiint_{V}\nab\cdot\vc{A}\,\dV=\oiint_{\partial V}\vc{A}\cdot\dS.
\end{equation}
\end{teorema}

\begin{teorema}[Ley de Gauss, forma diferencial: primera ecuación de Maxwell]\label{teo:gauss-dif}
En el vacío, el campo electrostático generado por una densidad volúmica de carga $\rho$ cumple
\begin{equation}\label{eq:gauss-dif}
  \boxed{\;\nab\cdot\vE(\vc{r})=\frac{\rho(\vc{r})}{\eo}\;}
\end{equation}
\end{teorema}
\begin{proof}
Sea $V$ un volumen cualquiera. Por la forma integral (\cref{eq:gauss-int}) y por el teorema de la divergencia
(\cref{eq:div-teorema}),
\[
  \iiint_{V}\nab\cdot\vE\,\dV=\oiint_{\partial V}\vE\cdot\dS=\frac{1}{\eo}\iiint_{V}\rho\,\dV .
\]
Igualando, $\iiint_{V}\bigl(\nab\cdot\vE-\rho/\eo\bigr)\dV=0$ para \emph{todo} volumen $V$. Si el integrando fuese no nulo en
algún punto, por continuidad tendría signo constante en un entorno suyo y la integral sobre ese entorno no sería cero;
por tanto, el integrando es idénticamente nulo.
\end{proof}

Esta ecuación pone de manifiesto que las \emph{fuentes} del campo eléctrico son las cargas, cuya densidad es $\rho$: donde
$\rho>0$ el campo «diverge» (nacen líneas de campo) y donde $\rho<0$ «converge». Es una consecuencia de la ley de Coulomb
y de su dependencia con el inverso del cuadrado de la distancia, y es válida siempre que estemos en el vacío. Falta aún
resolver el caso de las cargas puntuales: en el punto donde está la carga, $\rho$ diverge, y necesitamos una herramienta
matemática para escribir la ecuación. Es la delta de Dirac.

% -------------------------------------------------------------------------
\subsection{La delta de Dirac}\label{sec:delta}
% -------------------------------------------------------------------------
\begin{definicion}[Delta de Dirac]\label{def:delta}
La \emph{delta de Dirac} se define como el límite, cuando $\alpha\to\infty$, de una sucesión de funciones $f_\alpha(x)$ (por
ejemplo, un pulso rectangular de anchura $1/\alpha$ y altura $\alpha$, un triángulo de base $2/\alpha$ y altura $\alpha$ o una
gaussiana estrecha de área unidad) (\cref{fig:delta}):
\[
  \delta(x)=\lim_{\alpha\to\infty}f_\alpha(x),\qquad
  \delta(x)=\begin{cases}0,&x\neq0,\\\infty,&x=0,\end{cases}\qquad
  \int_{-\infty}^{\infty}\delta(x)\,\dif x=1 .
\]
Análogamente, $\delta(x-a)$ es nula salvo en $x=a$, donde diverge. Su propiedad esencial es la de \emph{muestreo}:
\begin{equation}\label{eq:delta-muestreo}
  \int_{-\infty}^{\infty}f(x)\,\delta(x-a)\,\dif x=f(a)
\end{equation}
para toda $f$ continua en $a$.
\end{definicion}

% Figura: sucesiones que definen la delta de Dirac.
\begin{figure}[ht]
\centering
\pgfplotsset{deltaax/.style={tema, width=0.335\textwidth, height=4.3cm, xmin=-1.3, xmax=1.3, ymin=0, ymax=4.4,
   xtick={-1,0,1}, ytick={0,1,2,4}, xlabel={$x$}, grid=major}}
\begin{tikzpicture}
  \begin{axis}[deltaax, name=A, ylabel={$f_\alpha(x)$}, title={\small\scshape (a)\; pulsos}]
    \addplot[azulclaro!80!azulacero, line width=1.1pt, const plot, mark=none] coordinates {(-1.3,0) (-0.5,0) (-0.5,1) (0.5,1) (0.5,0) (1.3,0)};
    \addplot[azulacero, line width=1.1pt, const plot, mark=none] coordinates {(-1.3,0) (-0.25,0) (-0.25,2) (0.25,2) (0.25,0) (1.3,0)};
    \addplot[azulprof, line width=1.1pt, const plot, mark=none] coordinates {(-1.3,0) (-0.125,0) (-0.125,4) (0.125,4) (0.125,0) (1.3,0)};
  \end{axis}
  \begin{axis}[deltaax, at={(A.east)}, anchor=west, xshift=0.6cm, yticklabels={}, title={\small\scshape (b)\; triángulos}]
    \addplot[azulclaro!80!azulacero, line width=1.1pt] coordinates {(-1.3,0) (-1,0) (0,1) (1,0) (1.3,0)};
    \addplot[azulacero, line width=1.1pt] coordinates {(-1.3,0) (-0.5,0) (0,2) (0.5,0) (1.3,0)};
    \addplot[azulprof, line width=1.1pt] coordinates {(-1.3,0) (-0.25,0) (0,4) (0.25,0) (1.3,0)};
  \end{axis}
  \begin{axis}[deltaax, at={(A.east)}, anchor=west, xshift=5.85cm, yticklabels={}, samples=200, domain=-1.3:1.3,
               title={\small\scshape (c)\; gaussianas}]
    \addplot[azulclaro!80!azulacero, line width=1.1pt] {1/(0.6*sqrt(2*pi))*exp(-x^2/(2*0.6^2))};
    \addplot[azulacero, line width=1.1pt] {1/(0.3*sqrt(2*pi))*exp(-x^2/(2*0.3^2))};
    \addplot[azulprof, line width=1.1pt] {1/(0.1*sqrt(2*pi))*exp(-x^2/(2*0.1^2))};
  \end{axis}
\end{tikzpicture}
\caption{La delta de Dirac como límite ($\alpha\to\infty$) de picos cada vez más estrechos y altos, todos de
\emph{área unidad}. (a)~Pulsos rectangulares de anchura $1/\alpha$ y altura $\alpha$ ($\alpha=1,2,4$).
(b)~Triángulos de base $2/\alpha$ y altura $\alpha$. (c)~Gaussianas $\frac{1}{\sigma\sqrt{2\pi}}\eul^{-x^{2}/2\sigma^{2}}$
con $\sigma=0{,}6;\,0{,}3;\,0{,}1$. Cuando el pico se afila, el área no cambia y el valor en $x=0$ diverge.}
\label{fig:delta}
\end{figure}


Estrictamente, la delta no es una función sino una \emph{función generalizada} o \emph{distribución}, y solo adquiere sentido
bajo el signo integral. La propiedad de muestreo se entiende así: como $\delta(x-a)=0$ para $x\neq a$, en la integral solo cuenta el
valor de $f$ en $x=a$, que sale fuera de ella: $\int f(x)\delta(x-a)\dif x=f(a)\int\delta(x-a)\dif x=f(a)$.

\begin{ejemplo}[Cálculos con la delta]\label{ej:delta}
Con \cref{eq:delta-muestreo} (la delta cae en el interior del intervalo de integración):
\[
  \int_{-\infty}^{\infty}(3x-1)\,\delta(x+1)\,\dif x=3(-1)-1=-4,\qquad
  \int_{0}^{3}x^{3}\,\delta(x-2)\,\dif x=2^{3}=8 .
\]
\end{ejemplo}

En tres dimensiones se define $\delta^{3}(\vc{r})=\delta(x)\,\delta(y)\,\delta(z)$ y, para un punto $\vc{a}$ del espacio,
\begin{equation}\label{eq:delta3}
\begin{gathered}
  \delta^{3}(\vc{r}-\vc{a})=\begin{cases}0,&\vc{r}\neq\vc{a},\\\infty,&\vc{r}=\vc{a},\end{cases}\qquad
  \int_{V}\delta^{3}(\vc{r}-\vc{a})\,\dV=\begin{cases}1,&\vc{a}\in V,\\0,&\vc{a}\notin V,\end{cases}\\[4pt]
  \int_{V}f(\vc{r})\,\delta^{3}(\vc{r}-\vc{a})\,\dV=\begin{cases}f(\vc{a}),&\vc{a}\in V,\\0,&\vc{a}\notin V.\end{cases}
\end{gathered}
\end{equation}

\begin{ejemplo}[Delta tridimensional]\label{ej:delta3}
$\displaystyle\int_{\R^{3}}\bigl(r^{2}+\vc{r}\cdot\vc{a}+a^{2}\bigr)\,\delta^{3}(\vc{r}-\vc{a})\,\dV=a^{2}+a^{2}+a^{2}=3a^{2}$.
\end{ejemplo}

Con ella podemos precisar qué ocurre con el campo de una carga puntual.

\begin{teorema}[Divergencia del campo coulombiano y laplaciano de $1/R$]\label{teo:delta}
Sea $\vc{R}=\vc{r}-\vc{r}\,'$. Entonces, derivando respecto de $\vc{r}$,
\begin{equation}\label{eq:div-delta}
  \nab\cdot\Bigl(\frac{\uR}{R^{2}}\Bigr)=\nab\cdot\Bigl(\frac{\vc{R}}{R^{3}}\Bigr)=4\pi\,\delta^{3}(\vc{R}),
  \qquad
  \Lap\Bigl(\frac{1}{R}\Bigr)=-4\pi\,\delta^{3}(\vc{R}).
\end{equation}
\end{teorema}
\begin{proof}
Tomemos el origen en la fuente ($\vc{r}\,'=\vc{0}$, $\vc{R}=\vc{r}$). Para $R\neq0$ el campo es radial, $A_R=1/R^{2}$, y en coordenadas
esféricas $\nab\cdot\vc{A}=\frac{1}{R^{2}}\pd{}{R}\bigl(R^{2}A_R\bigr)=\frac{1}{R^{2}}\pd{(1)}{R}=0$. Pero en $R=0$ el campo diverge. Para
ver qué ocurre allí, integramos sobre un volumen $V$ cualquiera con el teorema de la divergencia y el \cref{lema:omega-4pi}:
\[
  \int_{V}\nab\cdot\Bigl(\frac{\uR}{R^{2}}\Bigr)\dV=\oiint_{\partial V}\frac{\uR\cdot\dS}{R^{2}}=\oiint_{\partial V}\dOm
  =\begin{cases}0,&\text{si }R=0\text{ no está en }V,\\4\pi,&\text{si }R=0\text{ sí está en }V.\end{cases}
\]
Una cantidad que vale cero fuera del origen y cuya integral sobre cualquier volumen que lo contiene es $4\pi$ es, por definición,
$4\pi\,\delta^{3}(\vc{R})$. Para el laplaciano, como $\nab(1/R)=-\vc{R}/R^{3}$, se tiene
$\Lap(1/R)=\nab\cdot\nab(1/R)=-\nab\cdot(\vc{R}/R^{3})=-4\pi\delta^{3}(\vc{R})$.
\end{proof}

\begin{observacion}[Ni cero ni infinito]\label{obs:delta-div}
En coordenadas esféricas, $\Lap(1/R)=\frac{1}{R^{2}}\pd{}{R}\bigl(R^{2}\pd{}{R}\frac{1}{R}\bigr)$ vale cero para $R\neq0$ y no está definido en
$R=0$, donde diverge, igual que $\delta(R)$. El teorema precisa cuánto: lo suficiente para que la integral sobre un entorno
del origen sea $-4\pi$. Un cálculo ingenuo que olvide el origen diría que $\nab\cdot(\uR/R^{2})=0$ en todas partes y
que el flujo $4\pi$ de este campo sale «de la nada»; la delta es lo que restituye la coherencia entre la divergencia y el flujo.
\end{observacion}

\begin{corolario}[Ley de Gauss diferencial de una carga puntual]\label{cor:gauss-puntual}
Para una carga puntual $\qp$ en $\vc{r}\,'$, $\rho(\vc{r})=\qp\,\delta^{3}(\vc{r}-\vc{r}\,')$ y
\[
  \nab\cdot\vE(\vc{r})=\frac{\qp}{4\pi\eo}\,\nab\cdot\Bigl(\frac{\vc{R}}{R^{3}}\Bigr)=\frac{\qp}{4\pi\eo}\,4\pi\,\delta^{3}(\vc{R})
  =\frac{\qp}{\eo}\,\delta^{3}(\vc{R})=\frac{\rho}{\eo}.
\]
La divergencia del campo de una carga puntual es nula en todas partes salvo en el punto que ocupa la carga, donde diverge.
\end{corolario}

Con la delta podemos dar una segunda demostración, directa, de la ley de Gauss diferencial \emph{sin pasar por la forma integral}.

\begin{proof}[Demostración directa de \cref{eq:gauss-dif}]
Partimos del campo de una distribución volúmica (\cref{eq:campo-vol}) y derivamos respecto del punto campo $\vc{r}$, es decir, respecto de las
coordenadas del punto campo y no de las de la fuente, de modo que $\nab$ actúa solo sobre $\vc{R}/R^{3}$ y se puede introducir bajo
la integral:
\[
  \nab\cdot\vE(\vc{r})=\frac{1}{4\pi\eo}\int_{V'}\rho(\vc{r}\,')\,\nab\cdot\Bigl(\frac{\vc{R}}{R^{3}}\Bigr)\dVp
  =\frac{1}{4\pi\eo}\int_{V'}\rho(\vc{r}\,')\,4\pi\,\delta^{3}(\vc{r}-\vc{r}\,')\,\dVp=\frac{\rho(\vc{r})}{\eo},
\]
por el \cref{teo:delta} y la propiedad de muestreo (\cref{eq:delta3}).
\end{proof}

% -------------------------------------------------------------------------
\subsection{Cálculo del campo con la ley de Gauss: simetrías}\label{sec:gauss-simetria}
% -------------------------------------------------------------------------
Si la distribución de carga posee suficiente simetría, la ley de Gauss permite calcular $\vE$ de manera sencilla. Hay que hacer
hincapié en que la ley es válida \emph{siempre}, pero que despejar $\vE$ de ella solo es posible si existe la simetría adecuada:
hay que elegir una superficie de integración (\emph{superficie gaussiana}, una construcción matemática que no tiene por qué ser real)
en la cual, en todos sus puntos, $\vE$ sea constante en módulo y, además, paralelo o perpendicular a $\dS$. Sólo hay tres
simetrías que funcionan:

\begin{center}\small
\renewcommand{\arraystretch}{1.3}
\begin{tabular}{@{}lll@{}}
\toprule
Simetría & Superficie de Gauss & Campo\\
\midrule
Esférica & esfera de radio $r$ & $\vE=E(r)\,\ur$\\
Cilíndrica & cilindro coaxial de radio $r$ y altura $\ell$ & $\vE=E(r)\,\urho$\\
Plana & cilindro (o paralelepípedo) recto que atraviesa el plano & $\vE=E(z)\,\uz$\\
\bottomrule
\end{tabular}
\end{center}

En esos casos, $\oiint_{S}\vE\cdot\dS=E\,S$ y se despeja $E(\vc{r})=\Qenc(\vc{r})/(\eo\,S(\vc{r}))$. Los seminarios
\cref{sem:esfera,sem:hilo-gauss,sem:plano-gauss,sem:lamina,sem:cascara,sem:cilindro} aplican el método a cada una de ellas.

\begin{ejemplo}[De los campos a las fuentes]\label{ej:inverso}
La ley de Gauss diferencial también se usa en sentido inverso: dado el campo, \emph{deducir la carga}. Si $\vE(\vc{r})=k\,z\,\uz$,
entonces $\nab\cdot\vE=k$ y por \cref{eq:gauss-dif} la densidad es $\rho=k\,\eo$: una distribución volúmica de carga \emph{uniforme} en todo el espacio.
\end{ejemplo}
